\chapter{Methodology I: Simplification Approaches}
\label{sec:methodology-i}

\section{Comparison design}
\label{sec:comparison-design}

The comparison is built as a grid of three approaches against two granularities.
The approaches are the comparison simplifier, a \gls{seqseq} model fine-tuned for this project, and an open-weight instruction-tuned \gls{llm} steered by a prompt; the granularities are the sentence and the document.

To answer \cref{rq:comparison}, the trained model is compared with two checkpoints.
\begin{enumerate}
	\item The \textbf{untuned checkpoint} is the \gls{bart}-base model before fine-tuning.
	\item The \textbf{comparison simplifier} is the off-the-shelf third-party fine-tuned model described in \cref{sec:related-baselines}.
	      It is roughly three times as large and was trained on the same WikiLarge corpus.
\end{enumerate}

The fine-tuned model on sentence-level and the comparison simplifier therefore differ mainly in model size and in their fine-tuning runs, which were set up independently.
There is no document counterpart, since the comparison simplifier was fine-tuned on single sentences (\cref{sec:related-baselines}).

\section{Corpus assessment and preparation}
\label{sec:corpus-preparation}

\Cref{tab:corpora} specifies the five assessed corpora and the decision taken on each.

\begin{table}[!tb]
	\centering
	\footnotesize
	\setlength{\tabcolsep}{3pt}
	\caption{Corpus Assessment and Selection Decisions}
	\label{tab:corpora}
	\begin{tabular}{@{}P{1.9cm}P{3.2cm}P{4.3cm}P{2.6cm}@{}}
		\toprule
		\textbf{Corpus (unit)}            & \textbf{Rows as used here}                                                                                     & \textbf{Convention and alignment}                                                                                               & \textbf{Decision}                                                           \\
		\midrule
		WikiLarge\newline (sent.~pairs)   & \num{121657} / 410 / 121 filtered; \num{123862} / 417 / 121 raw; source \glstext{p90} \({\approx}52\)~tokens   & natural case; doubly-encoded characters in a minority of rows; automatic, aggregates earlier collections, among them \gls{pwkp} & \textbf{Adopted} for sentence-level training                                \\
		\addlinespace
		WikiSmall\newline (sent.~pairs)   & \({\approx}89\text{\gls{thousand}}\) pairs in the distribution inspected; splits unused                        & named-entity placeholders throughout (\texttt{PERSON@1}); automatic, sentence-level \gls{tfidf}                                 & \textbf{Rejected} on named-entity anonymisation                             \\
		\addlinespace
		D-Wikipedia\newline (doc.~pairs)  & \num{131739} train (filtered); \num{8000} test; mean 180, median 120, \glstext{p95} 510, max \num{1168} tokens & lowercased, Penn-Treebank pre-tokenised, one document per line; automatic                                                       & \textbf{Adopted} for document-level training                                \\
		\addlinespace
		SWiPE\newline (doc.~pairs)        & \({\approx}145\text{\gls{thousand}}\) pairs; splits unused                                                     & automatic, with human edit-category annotation                                                                                  & \textbf{Rejected}; built for classifying edit types, not for plain training \\
		\addlinespace
		ASSET\newline (sent., multi-ref.) & no train split; validation used for demonstrations; 359 test                                                   & natural case; human rewrites, ten per source sentence                                                                           & \textbf{Adopted} for sentence-level evaluation only                         \\
		\bottomrule
	\end{tabular}
	\note{
		All corpora are English only; counts describe the distributions used here, not the corpora as published: \textcite[589]{zhang-lapata-2017-sentence} report \num{296402} pairs for WikiLarge.
		Lengths are source lengths in tokens; D-Wikipedia's come from a sample of \num{1500} source documents, and its validation split and the two rejected corpora were not profiled.
		Sources, in table order: \parencites{zhang-lapata-2017-sentence}{zhu-etal-2010-monolingual}{sun-etal-2021-document}{laban-etal-2023-swipe}{alva-manchego-etal-2020-asset}.
	}
\end{table}

\paragraph{The WikiSmall distribution\protect\footnote{\hfdataset{cestwc/adapted-wikismall}} was rejected on inspected evidence.}
Throughout its text named entities are replaced by placeholders such as \texttt{PERSON@1} and \texttt{LOCATION@1}.
A model trained on such pairs would learn to produce placeholders where a page has names.
The placeholders appear to be a property of this distribution rather than of the original corpus: \textcite[1354]{zhu-etal-2010-monolingual} describe no named-entity replacement step for \gls{pwkp}, the corpus distributed as WikiSmall.
They most likely come from the release of \textcite[588--589]{zhang-lapata-2017-sentence}, who replaced named entities with tokens such as \texttt{PER@1} and whose WikiSmall training set has \num{89042} pairs, the size of the distribution inspected here.

\paragraph{SWiPE\protect\footnote{\ghrepo{salesforce/simplification}} was rejected as built for a different question.}
It pairs English and Simple English Wikipedia pages reconstructed from their revision histories, with a subset of about \num{5000} pairs annotated for 19 categories of edits \parencite[10674--10675]{laban-etal-2023-swipe}.
That annotation serves the classification of edit types, whereas flat document pairs were needed for the \gls{seqseq} training.

\pagebreak

\paragraph{A pre-cleaned WikiLarge mirror was adopted, but it required further cleaning.}
The mirror is the one introduced in \cref{sec:related-baselines}; its curator had already deduplicated it and filtered it by token length, compression ratio, and lexical similarity \parencite[3--4]{cohen2026simplify}.
One defect survived that filtering: mojibake, the doubly-encoded byte sequences produced when UTF-8 text is decoded as Latin-1 somewhere in an alignment pipeline.
Affected rows are dropped rather than repaired, since the standard repair tool, \texttt{ftfy}, does not recognise this pattern: the broken byte sequences are no longer contiguous.

Two filters exist (\cref{tab:corpus-rows}).
A strict filter drops any row containing a bare \texttt{\^{a}} or \texttt{\~{A}}, while a targeted one matches seven specific broken sequences.
The strict filter drops every row the targeted one drops, plus \num{4001} further training rows (\qty{3.3}{\percent}) that inspection identified as mojibake the targeted patterns miss.
A spot check found legitimate accented text surviving the strict filter.
The deployed sentence checkpoint was trained under the targeted filter, while the corpus statistics in the repository's data statement (\path{research/data/README.md} and \path{stats.json}) are those of the strict filter.

\begin{table}[htbp]
	\centering
	\small
	\caption{Comparison of WikiLarge Split Sizes Under the Two Mojibake Filters}
	\label{tab:corpus-rows}
	\begin{tabular}{@{}lrrr@{}}
		\toprule
		\textbf{Split} & \textbf{Raw}                          & \textbf{Strict filter}                                              & \textbf{Targeted filter}                                            \\
		\midrule
		train          & \({\approx}123\text{\gls{thousand}}\) & \({\approx}117\text{\gls{thousand}}\) \ (kept \qty{95.0}{\percent}) & \({\approx}121\text{\gls{thousand}}\) \ (kept \qty{98.2}{\percent}) \\
		validation     & 417                                   & 397 \ (kept \qty{95.2}{\percent})                                   & 410 \ (kept \qty{98.3}{\percent})                                   \\
		test           & 121                                   & 121 \ (kept \qty{100}{\percent})                                    & 121 \ (kept \qty{100}{\percent})                                    \\
		\midrule
		all splits     & \({\approx}124\text{\gls{thousand}}\) & \({\approx}118\text{\gls{thousand}}\) \ (kept \qty{95.0}{\percent}) & \({\approx}122\text{\gls{thousand}}\) \ (kept \qty{98.2}{\percent}) \\
		\bottomrule
	\end{tabular}
	\note{
		Percentages are of the raw split.
		Exact training-row counts are \num{123862} raw, \num{117656} strict and \num{121657} targeted.
	}
\end{table}

\paragraph{D-Wikipedia needed its own loader.}
It is distributed not through the \gls{hf} Hub but as raw files from the paper's own repository, with the training split in a compressed archive.
The cleaning code was reused unmodified, since it operates on whole strings regardless of whether a string is one sentence or one article.
The two granularities' corpora were therefore prepared identically, additionally they supply their pairs in the same flat source--target form.

\paragraph{Three scope tiers bound the cost of each run.}
The pipeline defines them at both granularities: \num{200} rows (\texttt{prove\_loop}), \num{20000} rows (\texttt{reduced}) and the entire cleaned split (\texttt{full}), so that every configuration runs on \num{200} rows before compute is spent on the full split.

\section{Fine-tuning sequence-to-sequence models}
\label{sec:finetuning}

Both fine-tuned checkpoints start from \gls{bart}-base \parencite[7872]{lewis-etal-2020-bart} and were trained with the \texttt{Seq2SeqTrainer} of the Transformers library \parencite{wolf-etal-2020-transformers}.
Their data collator sets \texttt{label\_pad\_token\_id=-100}, which excludes padding from the loss.
The two granularities share the entire pipeline (tokenisation, trainer setup, decoding, and significance testing), and a single flag selects the dataset structure.
Differences between the granularities' results should therefore come from the data, not from the implementation.

Early stopping after four evaluations without improvement was enabled, as was selection of the best checkpoint by validation loss.
The trainer evaluated every quarter epoch and saved a checkpoint once per epoch.
The sentence-level run shows the selection working, since epoch~4 was selected over epoch~5 (\num{0.4568} against \num{0.4590}).
Because evaluation and saving run at different intervals, the best validation loss the trainer records may fall at a step for which no checkpoint exists.
In that case, the figure quoted is the one the selected weights record.
Since \qty{19.02}{\percent} of the WikiLarge validation split also appears in the training data (\cref{sec:corpus-results}), the validation losses reported here are training diagnostics, not estimates of generalisation.
Checkpoint selection is unaffected, because every checkpoint is scored on the same validation split.

\paragraph{The length limits follow the training setup and the measured document lengths.}
The sentence models generate with \texttt{max\_length=64}, the value their fine-tuning used, which aligns inference with training at the cost of cutting off any output longer than 64~tokens.
The \gls{p90} of the source length is \({\approx}52\)~tokens.
The document models use an input cap of \num{512}~tokens, chosen from the measured token-length distribution of the training corpus under the available memory budget (\cref{fig:doclen}).
\Gls{bart}'s positional limit is \num{1024}~tokens.
\Cref{sec:granularity-findings} uses the same distribution to decide whether page sections are merged into one document input.

\begin{figure}[htbp]
	\centering
	\includegraphics[width=0.8\textwidth]{fig_doclen_distribution.pdf}
	\caption{Source-Document Length Distribution Relative to the \num{512}-Token Input Cap}
	\label{fig:doclen}
	\note{
		The six markers are measured quantiles; the curve joining them is a monotone interpolation, not measured data.
	}
\end{figure}

\paragraph{The training configuration depends on the device} (\cref{tab:training-config}).
On Apple Silicon, dynamic per-batch padding let reserved memory grow past \qty{13}{\glsxtrshort{gib}} within roughly a hundred steps, against \({\approx}\qty{4.5}{\glsxtrshort{gib}}\) with fixed-length padding.
Sentence training on Apple Silicon therefore uses fixed padding, gradient accumulation, and gradient checkpointing.
At document level, the progress estimate of a full-scope run exceeded the session cap of the hosted Google Colab tier \parencite{google2026colabfaq}, the full-scope run therefore was completed on an RTX~A5000 in under \qty{3}{\hour}.

\begin{table}[htbp]
	\centering
	\footnotesize
	\setlength{\tabcolsep}{4pt}
	\caption{Training Configuration by Granularity, Scope, and Device}
	\label{tab:training-config}
	\begin{tabular}{@{}lccccl@{}}
		\toprule
		\textbf{Granularity; scope}      & \textbf{Rows}                                           & \textbf{Epochs} & \textbf{Batch} & \textbf{max\_len} & \textbf{Device} \\
		\midrule
		sentence; \texttt{prove\_loop}   & 200                                                     & 1               & 8              & 64                & MPS             \\
		sentence; \texttt{reduced}       & \({\approx}20\text{\gls{thousand}}\)                    & 2               & 16             & 64                & MPS             \\
		\textbf{sentence; \texttt{full}} & \textbf{\boldmath\({\approx}121\text{\gls{thousand}}\)} & \textbf{5}      & \textbf{16}    & \textbf{64}       & \textbf{CUDA}   \\
		document; \texttt{reduced}       & \({\approx}20\text{\gls{thousand}}\)                    & 2               & 8              & 512               & CUDA            \\
		\textbf{document; \texttt{full}} & \textbf{\boldmath\({\approx}131\text{\gls{thousand}}\)} & \textbf{5}      & \textbf{32}    & \textbf{512}      & \textbf{RTX}    \\
		\bottomrule
	\end{tabular}
	\note{
		Bold marks the two deployed checkpoints (exact row counts: \Cref{tab:corpus-rows}).
		Devices are \gls{mps} on Apple Silicon, CUDA on a hosted Google Colab T4 in \gls{fp16}, and RTX on an RTX~A5000 in \gls{bf16}.
		All runs share a learning rate of \num{3e-5}, weight decay~\num{0.01}, seed~42, early stopping with patience~4, and best checkpoint by validation loss.
	}
\end{table}

\paragraph{Decoding is by beam search, and the checkpoints are published.}
All reported \gls{seqseq} numbers use beam search (\texttt{num\_beams=4}, \texttt{no\_repeat\_ngram\_size=3}, \texttt{repetition\_penalty=1.2}).
Beam search explores several candidate continuations in parallel instead of committing to the locally most probable token at each step \parencite[224--225]{tunstall2023natural}.
One early run was logged as \emph{greedy} but passed no decoding arguments, so it used the checkpoint's own four-beam configuration.
A true greedy run was added afterwards, and the three runs span \num{0.26}~\gls{sari} (\cref{sec:sentence-results}).

Three checkpoints are published on the \gls{hf} Hub, with model cards recording configuration, metrics, and limitations.
At sentence level, the published checkpoint\footnote{\hfmodel{yunvs/bart-base-wikilarge-simplification}} is the deployed sentence checkpoint.
At document level, both the reduced-scope\footnote{\hfmodel{yunvs/bart-base-dwikipedia-simplification}} and the full-scope\footnote{\hfmodel{yunvs/bart-base-dwikipedia-simplification-full}} checkpoints are published.
The full-scope checkpoint is the one the document results report and the backend serves.

\section{Prompted open-weight language models}
\label{sec:method-three}

\paragraph{Only open-weight \glspl{llm} were considered.}
A hosted \gls{api} would void the local-processing constraint of \cref{sec:architecture}.
The constraint has a known cost, since \textcite[13291]{kew-etal-2023-bless} report that closed-weight models provide significant gains over open-weight ones under few-shot prompting.
A local Ollama server \parencite{ollama-docs-api} serves the model in 4-bit quantised GGUF form, a model file format Ollama imports \parencite{ollama-docs-import}.
At roughly \qty{4.7}{\glsxtrshort{gb}}, the quantised 7\gls{billion} model fits alongside the \gls{bart} checkpoints on a \qty{16}{\glsxtrshort{gb}} machine, which the roughly \qty{14}{\glsxtrshort{gb}} half-precision version would not.
Two instruction-tuned models of the Qwen2.5 family \parencite[1]{qwen2024qwen25} were evaluated, one with 7\gls{billion}\footnote{\ollamamodel{qwen2.5:7b-instruct-q4_K_M}} and one with 3\gls{billion}\footnote{\ollamamodel{qwen2.5:3b-instruct-q4_K_M}} parameters.
The backend sends the prompted model one request at a time, so the prompted approach runs at a batch size of one.
It does not use the batching the \gls{seqseq} checkpoints rely on, which may have limited its throughput (\cref{sec:runtime}).

\paragraph{The sentence prompt is \gls{bless} Prompt~2 with an audience slot.}
It reproduces the prompt verbatim \parencite[13294]{kew-etal-2023-bless}, with one change: the hard-coded audience \emph{non-native speakers of English} becomes a slot that holds one of five audience values, such as \emph{children aged 8 to 12}.
The original phrase is the default.
\Cref{app:prompts} gives the exact template, the audience values, and the decoding settings.

ASSET's references were written under instructions for non-native speakers of English, the instructions \gls{bless}'s Prompt~2 repurposes \parencite[13294]{kew-etal-2023-bless}.
Pairing them with an instruction to simplify for children would set the demonstrations against the instruction.
Each audience therefore has its own demonstration pool: ASSET pairs for the default audience and hand-authored pairs for the other four (\cref{app:demos}).

\paragraph{The default demonstrations are fixed and pass four selection rules.}
\Gls{bless} samples its demonstrations at random and leaves the study of in-context learning strategies open (\cref{sec:related-baselines}).
Here, the same three demonstrations are used in every run: ASSET validation indices 285, 1516, and 1116 (\cref{app:demos}).
They were selected by four rules.
The first is a random draw at seed~42; the other three are filters, each added after an observed defect:
\begin{itemize}
	\item the source has at least 20~words, so that at least one demonstration can show a sentence split;
	\item the reference is well-formed, which excludes pairs with crowdworker punctuation typos;
	\item at least half of the reference's content words occur in the source.
\end{itemize}
The last filter excluded an ASSET reference that is fluent and well punctuated but has a different meaning from its source, the kind of misalignment \textcite[285]{xu-etal-2015-problems} found in \qty{17}{\percent} of \num{200} sampled \gls{pwkp} pairs.
Together, the three demonstrations show all three operations the instruction names (paraphrasing, compression, and splitting).
None of the three demonstrations occurs in the WikiLarge training data.

\paragraph{Failed outputs are classified by reason code.}
The output check of the prompted approach, one of the guards described in \cref{sec:guards}, returns a \emph{reason code} instead of a pass or fail.
Some codes cover failures typical of a prompted \gls{llm}, such as prompt echo and refusal.
The others cover failures any generator can produce: empty output, output without words, degenerate repetition, over-long output, unrelated short output, timeout, and error.

The counts per code feed the error categorisation directly.
When the check fails, the original text is kept, so a bad output leaves the sentence unsimplified.
Output identical to the input does not count as a failure: copying the source is a documented behaviour on this task \parencite[13298]{kew-etal-2023-bless}.
Such output is therefore measured and reported per condition as the unchanged rate.

\pagebreak

\paragraph{The document prompt adapts the sentence prompt.}
\Gls{bless} is sentence-only, so what carries over is the prompt's structure.
\Cref{app:prompts} prints the adapted template beside the original.
Three differences are recorded: merging and reordering are added to the permitted operations; the output is constrained to continuous prose with no preamble, headings, or bullet points; and the frame is \texttt{Text:}/\texttt{Simplified text:} instead of \texttt{Complex:}/\texttt{Simple:}.
The document prompt carries no demonstrations by default: each would be a whole document pair, which would make the prompt much longer.
Such pairs would moreover come from D-Wikipedia, which is lowercased and pre-tokenised, unlike typical web-page text.
The prompted approach therefore does not denote the same configuration at the two granularities.

Two decoding configurations are likewise reported separately: greedy decoding at temperature~0 for serving, so that re-simplifying a page is stable, and \gls{bless}'s own sampling settings \parencite[13293]{kew-etal-2023-bless} for evaluation, so that the numbers remain comparable to the paper.
The extension and the results table therefore do not describe the same system.

\section{Evaluation protocol}
\label{sec:evaluation-protocol}

\paragraph{Each granularity is scored on its own benchmark, but by the same evaluation code.}
Sentence-level conditions are scored on the ASSET test split \parencite[4670--4671]{alva-manchego-etal-2020-asset}.
Document-level conditions are scored on D-Wikipedia's own test split.

Scoring runs in a standalone evaluation pipeline, not in the metric code built into the trainer.
\Gls{sari} and \gls{fkgl} are computed with EASSE, a toolkit for evaluating sentence simplification \parencite[49--50]{alva-manchego-etal-2019-easse}.
Every approach passes through the same data loader, the same reference set per benchmark, and the same metric calls, so differences between approaches cannot come from the scoring code.

\paragraph{Significance is tested by paired bootstrap resampling.}
Improvements in \gls{sari} and \gls{dsari} are tested \parencite[393]{koehn-2004-statistical} with \num{1000} resamples, applying the same resampled indices to both systems so that pairing is preserved; this provides the empirical basis for refuting Hypothesis~\cref{hyp:finetuning}.
Where the fine-tuned system wins every resample, the bootstrap estimate is \(\gls{pval} = 0\).
This estimate is reported as \(\gls{pval} < 0.001\), because \num{1000} resamples cannot distinguish a truly zero probability from a very small one.
\Gls{lens} and BERTScore are reported descriptively; testing their differences for significance is outside the scope of this thesis.

\pagebreak

\paragraph{Metric preprocessing is part of the metric.}
Where a metric carries structural penalties, its input preprocessing changes what it measures.
Normalising into D-Wikipedia's convention has no effect on \gls{sari} and \gls{fkgl}, which re-tokenise their input anyway, and changes \gls{dsari}, whose sentence-count penalty reads exactly the boundaries that normalisation creates or destroys.
\Cref{sec:document-results} measures this by rescoring cached generations under two versions of the normaliser.

\paragraph{Deployment inference and benchmark inference run in different places.}
Deployment inference runs locally, because privacy is a design constraint.
Benchmark inference runs on an RTX~A5000 \gls{gpu} instead.
ASSET and D-Wikipedia are public datasets that contain no user data, so the privacy constraint does not apply to them.
Local hardware was also too slow for the full five-condition evaluation, which the \gls{gpu} completed in 42~minutes.

\paragraph{Provenance is recorded, and the \gls{dsari} reimplementation is tested.}
From the first run on the RTX~A5000 onwards, a staged runner committed to the repository produced every number on that machine, at one code revision.
Before it spends any \gls{gpu} time, the runner executes a correctness check, which caught a metric-normalisation defect before a training run.
Every artefact records the software environment that produced it, including the exact commit of the \gls{sari} implementation.

A differential test compares the \gls{dsari} reimplementation with the reference implementation: across 66 cases drawn from real D-Wikipedia documents, the two must agree to full floating-point precision.
Nine further invariant checks run offline.
Since the reference implementation ships no tests of its own, this test is what makes the document scores comparable with published \gls{dsari} figures.
The reimplementation keeps two properties of the reference implementation, so that it measures \gls{dsari} as published.
Deletion contributes a precision term instead of an F-score, as in the definition of \gls{sari} \parencite[405]{xu-etal-2016-optimizing}, and the reference implementation computes one intermediate delete-recall term that it never uses.
